\section{Modules and module homomorphisms}
\label{am-ch2-modules}
\begin{definition}[Module]
\label{am-def-module}
\source{atiyah-macdonald-commutative-algebra:page-0026}
An $A$-module is an abelian group $M$ with an action of $A$ satisfying
$a(x+y)=ax+ay$, $(a+b)x=ax+bx$, $(ab)x=a(bx)$, and $1x=x$.  An
$A$-linear map $f:M\to N$ preserves addition and scalar multiplication.
The set $\Hom_A(M,N)$ is an $A$-module by pointwise operations; pre- and
post-composition induce the usual maps on Hom.  Evaluation at $1$ gives a
natural isomorphism $\Hom_A(A,M)\simeq M$.
\end{definition}
\begin{definition}[Submodule, quotient, kernel, image]
\label{am-def-submodule-quotient}
\source{atiyah-macdonald-commutative-algebra:page-0027}
A submodule is an additive subgroup stable under $A$.  If $M'\subseteq M$,
$M/M'$ is an $A$-module.  For $f:M\to N$, $\Ker f$ and $\Im f$ are
submodules and $\Coker f=N/\Im f$.
\end{definition}
\begin{proposition}[First isomorphism theorem for modules]
\label{am-prop-2-1a}
\source{atiyah-macdonald-commutative-algebra:page-0028}
If $M'\subseteq\Ker f$, the induced map $M/M'\to N$ has kernel
$\Ker f/M'$.  In particular $M/\Ker f\simeq\Im f$.
\uses{am-def-submodule-quotient}
\end{proposition}
\begin{proof}
The induced map sends the class of $x$ to $f(x)$; this is well-defined because
$f(M')=0$.  Its kernel consists exactly of classes represented by elements of
$\Ker f$, and the last assertion is the first isomorphism theorem.
\end{proof}
\begin{definition}[Annihilator and faithfulness]
\label{am-def-ann-module}
\source{atiyah-macdonald-commutative-algebra:page-0028}
For submodules $N,P\subseteq M$, $(N:P)=\{a\in A:aP\subseteq N\}$ and
$\Ann(M)=(0:M)$.  If $\mathfrak a\subseteq\Ann(M)$, the action factors through
$A/\mathfrak a$.  The module is faithful when $\Ann(M)=0$.
\end{definition}
\begin{proposition}[Submodule quotients]
\label{am-prop-2-1}
\dcref{2.1}
\source{atiyah-macdonald-commutative-algebra:page-0028}
For $L\supseteq M\supseteq N$,
$(L/N)/(M/N)\simeq L/M$; and for $M_1,M_2\subseteq M$,
$(M_1+M_2)/M_1\simeq M_2/(M_1\cap M_2)$.
\end{proposition}
\begin{proof}
The first map is $x+N\mapsto x+M$, with kernel $M/N$.  The composite
$M_2\to M_1+M_2\to(M_1+M_2)/M_1$ is onto with kernel $M_1\cap M_2$.
\end{proof}

\section{Operations, sums, and products}
\label{am-ch2-operations}
\begin{remark}[Lattice operations]
\label{am-rem-module-lattice}
\source{atiyah-macdonald-commutative-algebra:page-0028}
Sums (finite sums for an arbitrary family) and intersections of submodules are
submodules and form a complete lattice.  For an ideal $\mathfrak a$,
$\mathfrak aM$ is generated by finite sums $\sum a_ix_i$; the quotient
$M/N$ is naturally an $A/\mathfrak a$-module when $\mathfrak aM\subseteq N$.
\end{remark}
\begin{definition}[Generators and finite generation]
\label{am-def-module-generators}
\source{atiyah-macdonald-commutative-algebra:page-0029}
The submodule generated by $x$ is $Ax$.  A family generates $M$ if every
element is a finite $A$-linear combination of its members; $M$ is finitely
generated if a finite family does so.
\end{definition}
\begin{definition}[Direct sum, product, free module]
\label{am-def-direct-free}
\source{atiyah-macdonald-commutative-algebra:page-0029}
The direct sum $\bigoplus_iM_i$ consists of families with finite support, while
$\prod_iM_i$ has no support restriction.  A free module is a direct sum of
copies of $A$; $A^n$ denotes the finite free module with $n$ summands.
\end{definition}
\begin{proposition}[Finite generation and free quotients]
\label{am-prop-2-3}
\dcref{2.3}
\source{atiyah-macdonald-commutative-algebra:page-0030}
An $A$-module $M$ is finitely generated iff it is a quotient of $A^n$ for some
$n$.
\end{proposition}
\begin{proof}
Given generators $x_1,\ldots,x_n$, the map $A^n\to M$ sending the $i$th basis
vector to $x_i$ is onto, and the first isomorphism theorem applies.  Conversely,
images of the standard basis generate any quotient.
\end{proof}
\begin{proposition}[Determinant trick]
\label{am-prop-2-4}
\dcref{2.4}
\source{atiyah-macdonald-commutative-algebra:page-0030}
If $M$ is finitely generated, $\mathfrak a$ an ideal, and an endomorphism
$\varphi$ satisfies $\varphi(M)\subseteq\mathfrak aM$, then
\[\varphi^n+a_1\varphi^{n-1}+\cdots+a_n=0\]
as endomorphisms for some $a_i\in\mathfrak a$.
\end{proposition}
\begin{proof}
For generators $x_i$, write $\varphi(x_i)=\sum_j a_{ij}x_j$ with
$a_{ij}\in\mathfrak a$.  The matrix $(\delta_{ij}\varphi-a_{ij})$ annihilates
the column $(x_j)$; multiplying by its adjugate gives that its determinant
annihilates all generators.  Expanding the determinant gives the equation.
\end{proof}
\begin{corollary}[A scalar annihilating a module]
\label{am-cor-2-5}
\dcref{2.5}
\source{atiyah-macdonald-commutative-algebra:page-0030}
If $\mathfrak aM=M$ for a finitely generated $M$, some $x\equiv1\pmod{\mathfrak a}$
satisfies $xM=0$.
\uses{am-prop-2-4}
\end{corollary}
\begin{proof}
Apply Proposition~\ref{am-prop-2-4} to $\varphi=1_M$ and take
$x=1+a_1+\cdots+a_n$.
\end{proof}
\begin{proposition}[Nakayama's lemma]
\label{am-prop-2-6}
\dcref{2.6}
\source{atiyah-macdonald-commutative-algebra:page-0030}
If $M$ is finitely generated and $\mathfrak a\subseteq J(A)$, then
$\mathfrak aM=M$ implies $M=0$.
\end{proposition}
\begin{proof}
Corollary~\ref{am-cor-2-5} gives $x\equiv1\pmod{J(A)}$ with $xM=0$.
Proposition~\ref{am-prop-1-9} makes $x$ a unit, hence $M=0$.  Alternatively,
a minimal generating set would allow its last generator to be eliminated since
$1-a_n$ is a unit.
\uses{am-cor-2-5,am-prop-1-9,am-def-jacobson}
\end{proof}
\begin{corollary}[Nakayama for a quotient]
\label{am-cor-2-7}
\dcref{2.7}
\source{atiyah-macdonald-commutative-algebra:page-0031}
If $N\subseteq M$ are finitely generated, $\mathfrak a\subseteq J(A)$, and
$M=\mathfrak aM+N$, then $M=N$.
\uses{am-prop-2-6,am-def-submodule-quotient}
\end{corollary}
\begin{proof}
Apply Nakayama to $M/N$, for which $\mathfrak a(M/N)=M/N$.
\end{proof}
\begin{proposition}[Generators over a local ring]
\label{am-prop-2-8}
\dcref{2.8}
\source{atiyah-macdonald-commutative-algebra:page-0031}
If $(A,\mathfrak m)$ is local and the images of $x_1,\ldots,x_n$ form a basis
of $M/\mathfrak mM$ over $A/\mathfrak m$, then the $x_i$ generate $M$.
\uses{am-cor-2-7,am-def-local-ring}
\end{proposition}
\begin{proof}
Their span maps onto $M/\mathfrak mM$, so $M=N+\mathfrak mM$; use
Corollary~\ref{am-cor-2-7}.
\end{proof}

\section{Exact sequences}
\label{am-ch2-exact}
\begin{definition}[Exactness]
\label{am-def-exact}
\source{atiyah-macdonald-commutative-algebra:page-0031}
A sequence is exact at a term when the image of the preceding map equals the
kernel of the following map.  Thus $0\to M'\to M$ means injective,
$M\to M''\to0$ means surjective, and a short exact sequence
$0\to M'\to M\to M''\to0$ identifies $M/M'$ with $M''$.
\end{definition}
\begin{proposition}[Hom detects exactness]
\label{am-prop-2-9}
\dcref{2.9}
\source{atiyah-macdonald-commutative-algebra:page-0031}
For $M'\xrightarrow{u}M\xrightarrow{v}M''\to0$, exactness is equivalent to
exactness of $0\to\Hom(M'',N)\to\Hom(M,N)\to\Hom(M',N)$ for every $N$.
Dually, exactness of $0\to N'\to N\to N''$ is equivalent to exactness of
$0\to\Hom(M,N')\to\Hom(M,N)\to\Hom(M,N'')$ for every $M$.
\end{proposition}
\begin{proof}
The forward implications are composition.  Conversely, injectivity for all
$N$ makes $v$ surjective.  Taking $N=M''$ and the identity shows $vu=0$.
Taking $N=M/\Im u$ and its quotient map forces every element of $\Ker v$ into
the image of $u$.  The dual argument is the same with arrows reversed.
\end{proof}
\begin{proposition}[Snake lemma]
\label{am-prop-2-10}
\dcref{2.10}
\source{atiyah-macdonald-commutative-algebra:page-0032}
Given a commutative diagram with exact rows
\[
0\to M'\to M\to M''\to0\\
0\to N'\to N\to N''\to0,
\]
there is an exact sequence
\[
0\to\Ker f'\to\Ker f\to\Ker f''\xrightarrow{\delta}
\Coker f'\to\Coker f\to\Coker f''\to0.
\]
\end{proposition}
\begin{proof}
For $x''\in\Ker f''$, lift it to $x\in M$.  Then $f(x)$ lies in the image of
$N'$, say $f(x)=u'(y')$, and define $\delta(x'')$ as the class of $y'$.
Changing lifts changes $y'$ by an element of $\Im f'$, so $\delta$ is
well-defined.  Direct lifting and kernel computations prove exactness at each
term.
\end{proof}
\begin{definition}[Additive function]
\label{am-def-additive-function}
\source{atiyah-macdonald-commutative-algebra:page-0032}
A function $\lambda$ on a class of modules is additive if
$\lambda(M')-\lambda(M)+\lambda(M'')=0$ on every short exact sequence in the
class.
\end{definition}
\begin{proposition}[Alternating sum]
\label{am-prop-2-11}
\dcref{2.11}
\source{atiyah-macdonald-commutative-algebra:page-0033}
For an exact finite sequence $0\to M_0\to\cdots\to M_n\to0$ whose terms and
kernels lie in the class of an additive function $\lambda$,
$\sum_{i=0}^n(-1)^i\lambda(M_i)=0$.
\end{proposition}
\begin{proof}
Let $N_i$ be the image at $M_i$, with $N_0=N_{n+1}=0$.  The short exact
sequences $0\to N_i\to M_i\to N_{i+1}\to0$ give
$\lambda(M_i)=\lambda(N_i)+\lambda(N_{i+1})$; the alternating sum cancels.
\end{proof}

\section{Tensor products}
\label{am-ch2-tensor}
\begin{definition}[Tensor product]
\label{am-def-tensor}
\source{atiyah-macdonald-commutative-algebra:page-0033}
An $A$-bilinear map is linear in each variable.  The tensor product $M\otimes_A N$
is characterized by a universal bilinear map $(x,y)\mapsto x\otimes y$:
every bilinear map to $P$ factors uniquely through an $A$-linear map
$M\otimes_A N\to P$.
\end{definition}
\begin{proposition}[Existence and uniqueness of tensor products]
\label{am-prop-2-12}
\dcref{2.12}
\source{atiyah-macdonald-commutative-algebra:page-0033}
The tensor product exists and is unique up to a unique isomorphism compatible
with pure tensors.
\end{proposition}
\begin{proof}
Take the free module on $M\times N$ and quotient by the submodule generated
by additivity in each variable and the relations
$(ax,y)-a(x,y)$ and $(x,ay)-a(x,y)$.  The images $x\otimes y$ satisfy the
bilinear identities and generate.  A bilinear map kills the defining relations,
so descends uniquely.  Two universal pairs induce inverse maps by uniqueness.
\end{proof}
\begin{proposition}[Multilinear tensor products]
\label{am-prop-2-12-star}
\dcref{2.12*}
\source{atiyah-macdonald-commutative-algebra:page-0035}
For finitely many modules there is a universal multilinear tensor product
$M_1\otimes\cdots\otimes M_r$, unique up to unique compatible isomorphism.
\uses{am-def-tensor,am-prop-2-12}
\end{proposition}
\begin{proof}
Use the same free-module quotient construction with multilinearity relations;
the universal factorization and uniqueness proof are unchanged.
\end{proof}
\begin{corollary}[Finite support of tensor relations]
\label{am-cor-2-13}
\dcref{2.13}
\source{atiyah-macdonald-commutative-algebra:page-0034}
If $\sum_i x_i\otimes y_i=0$ in $M\otimes N$, then it already vanishes in
$M_0\otimes N_0$ for finitely generated submodules $M_0\subseteq M$ and
$N_0\subseteq N$.
\uses{am-prop-2-12}
\end{corollary}
\begin{proof}
In the free-module construction the relation is a finite sum of defining
relations.  Include in $M_0,N_0$ the finitely many coordinates appearing in
that sum.
\end{proof}
\begin{proposition}[Canonical tensor isomorphisms]
\label{am-prop-2-14}
\dcref{2.14}
\source{atiyah-macdonald-commutative-algebra:page-0035}
There are natural isomorphisms $M\otimes N\simeq N\otimes M$,
$(M\otimes N)\otimes P\simeq M\otimes(N\otimes P)$,
$(M\oplus N)\otimes P\simeq(M\otimes P)\oplus(N\otimes P)$, and
$A\otimes M\simeq M$, induced respectively by swapping, reassociating,
distributing pure tensors, and $a\otimes x\mapsto ax$.
\end{proposition}
\begin{proof}
Each displayed rule is bilinear or multilinear, hence induces a homomorphism
by the universal property.  Construct the inverse by the reverse rule; both
composites agree with the identity on pure tensors, which generate.
\end{proof}

\section{Restriction and extension of scalars}
\label{am-ch2-scalars}
\begin{definition}[Restriction and extension]
\label{am-def-scalars}
\source{atiyah-macdonald-commutative-algebra:page-0036}
For $f:A\to B$, a $B$-module is an $A$-module by restriction of scalars.  For
an $A$-module $M$, $M_B=B\otimes_A M$ has the $B$-action
$b(b'\otimes x)=bb'\otimes x$ and is extension of scalars.
\end{definition}
\begin{proposition}[Finite generation under scalars]
\label{am-prop-2-16}
\dcref{2.16}
\source{atiyah-macdonald-commutative-algebra:page-0037}
If $N$ is finitely generated over $B$ and $B$ finitely generated over $A$, then
$N$ is finitely generated over $A$.
\end{proposition}
\begin{proof}
Products of a finite set of $A$-generators of $B$ with a finite set of
$B$-generators of $N$ generate $N$ over $A$.
\end{proof}
\begin{proposition}[Finite generation after extension]
\label{am-prop-2-17}
\dcref{2.17}
\source{atiyah-macdonald-commutative-algebra:page-0037}
If $M$ is finitely generated over $A$, then $B\otimes_A M$ is finitely generated
over $B$.
\uses{am-def-scalars}
\end{proposition}
\begin{proof}
The tensors $1\otimes x_i$ generate whenever the $x_i$ generate $M$.
\end{proof}

\section{Exactness of tensoring and algebras}
\label{am-ch2-flat}
\begin{proposition}[Tensoring is right exact]
\label{am-prop-2-18}
\dcref{2.18}
\source{atiyah-macdonald-commutative-algebra:page-0037}
If $M'\to M\to M''\to0$ is exact, then
$M'\otimes N\to M\otimes N\to M''\otimes N\to0$ is exact.
\end{proposition}
\begin{proof}
The natural adjunction
$\Hom(M\otimes N,P)\simeq\Hom(M,\Hom(N,P))$ and Proposition~\ref{am-prop-2-9}
show exactness after applying $\Hom(-,P)$ for every $P$, hence exactness.
\uses{am-prop-2-9,am-def-tensor}
\end{proof}
\begin{definition}[Flat module]
\label{am-def-flat}
\source{atiyah-macdonald-commutative-algebra:page-0038}
An $A$-module $N$ is flat if $M\mapsto M\otimes_A N$ preserves every exact
sequence.
\end{definition}
\begin{proposition}[Criteria for flatness]
\label{am-prop-2-19}
\dcref{2.19}
\source{atiyah-macdonald-commutative-algebra:page-0038}
For an $A$-module $N$ the following are equivalent: (i) $N$ is flat; (ii)
tensoring every short exact sequence remains exact; (iii) $f\otimes1_N$ is
injective for every injective $f$; (iv) the same holds for injective maps of
finitely generated modules.
\end{proposition}
\begin{proof}
Split any long exact sequence into short exact sequences for (i)$\Leftrightarrow$(ii);
(ii)$\Leftrightarrow$(iii) uses right exactness.  (iii)$\Rightarrow$(iv) is
immediate.  For the converse, an element in the kernel of $f\otimes1$ is a
finite sum of pure tensors; by Corollary~\ref{am-cor-2-13} it lies in tensoring
of finitely generated submodules, where injectivity is assumed, so it is zero.
\uses{am-prop-2-18,am-cor-2-13}
\end{proof}
\begin{definition}[Algebra and finite type]
\label{am-def-algebra}
\source{atiyah-macdonald-commutative-algebra:page-0039}
An $A$-algebra is a ring $B$ with a homomorphism $A\to B$.  It is finite if
$B$ is finitely generated as an $A$-module, and of finite type if finitely many
elements generate it as an $A$-algebra, equivalently if a polynomial ring over
$A$ maps onto it.  Every ring is a $\Z$-algebra.
\end{definition}
\begin{remark}[Tensor product of algebras]
\label{am-rem-tensor-algebra}
\source{atiyah-macdonald-commutative-algebra:page-0039}
If $B,C$ are $A$-algebras, $B\otimes_A C$ has the unique multiplication
$(b\otimes c)(b'\otimes c')=bb'\otimes cc'$ and identity $1\otimes1$ making
it an $A$-algebra; the construction follows from multilinearity and the
universal property in Proposition~\ref{am-prop-2-12-star}.
\uses{am-prop-2-12-star,am-def-algebra}
\end{remark}
